\chapter{Mean Polytopes and Maximum Entropy Distributions}

\abstract{
    Beyond elementary: Study generic maximum entropy distributions.
    Motivation from the chapter before
}

\sect{Polytopes of Mean Parameters}

We in this section investigate properties of probability distributions based on their mean parameters.
Given a statistic $\sstat$, we first define a mean parameter to any distribution by the expectation of the statistic.

\begin{definition}
    \label{def:meanPolytope}
    Let there be a statistic $\sstat$ and a boolean base measure $\basemeasurewith$.
    We call the tensor
    \begin{align*}
        \meanparamwith
        = \contractionof{\probat{\shortcatvariables},\sencsstatat{\shortcatvariables,\selvariable}}{\selvariable}
    \end{align*}
    the mean parameter tensor to a distribution $\probat{\shortcatvariables}$.
    The set
    \begin{align*}
        \genmeanset
        = \left\{\contractionof{\probtensor,\sencsstat,\basemeasure}{\selvariable}\wcols\probwith\in\bmrealprobof{\basemeasure} \right\}
    \end{align*}
    is called the polytope of realizable mean parameters.
    Here we denote by $\bmrealprobof{\basemeasure}$ the set of all probability distributions representable with respect to $\basemeasure$ (see \defref{def:representationBaseMeasure}).
\end{definition}

% Convex Hull Characterization Polytope
While introduced here as a property of a distribution, the mean parameters will be central to probabilistic inference in \charef{cha:probReasoning}.
We in the reminder of this section prepare for this application and derive tensor network representations for distributions having sufficient statistics $\sstat$, depending on their corresponding mean parameter in the polytope.


\subsect{Representation by Convex Hulls}\label{sec:meanPolytopeConvexHull}

\begin{figure}[t]
    \begin{center}
        \input{tikz-pics/meanset_sketch_generic.tex}
    \end{center}
    \caption{Sketch of the mean polytope ${\genmeanset}$ to a statistic $\sstat$, which is minimal with respect to $\basemeasure$.
    The mean polytope is a bounded subset of $\parspace$ (here sketched as a 2-dimensional projection). %, where each mean parameter is one of the three cases $\meanparamof{1},\meanparamof{2}$ or $\meanparamof{3}$.
    For any vertex $\meanparamof{1}$ we find $\shortcatindices$ such that $\meanparamofat{1}{\selvariable}=\sencsstatat{\indexedshortcatvariables,\selvariable}$.
    Generic mean parameters $\meanparamof{2}$ outside the interior are on a face $\genfacesetof{\facecondset}$. %not reproduced by a member of the exponential family $\expfamily$, but by a member of the family $\expfamilyof{\sstat,\secbasemeasure}$ where $\secbasemeasure$ is a refined base measure deterimined by \algoref{alg:baseMeasureRefinement}.
%are reproducible by \HybridLogicNetworks{} with statistic $\hlnstat$ and refined base measure $\secbasemeasure$.
    Interior points $\meanparamof{3}\in\interiorof{\hlnmeanset}$ are exactly those reproducible by positive distributions with respect to $\basemeasure$.  %member of $\expfamily$, which canonical parameter is found by the backward map.
    }\label{fig:meansetSketchGeneric}
\end{figure}

First of all, we provide a simple characterization of the sets of mean parameters as the convex hull of the slices to the selection encoding of the statistic (see \figref{fig:meansetSketchGeneric}).
Convex hulls of finite vectors are called $\mathcal{V}$-polytopes (see Lecture~1 in \cite{ziegler_lectures_2013}).

\begin{theorem}
    \label{the:meanPolytopeConvHull}
    For any statistic $\sstat$ the polytope of mean parameters is the convex hull of the slices of $\sencsstat$ with fixed indices to $\shortcatvariables$, that is
    \begin{align*}
        \genmeanset
        = \convhullof{\sencsstatat{\indexedshortcatvariables,\selvariable}\wcols\shortcatindices\in\facstates\ncond\basemeasureat{\indexedshortcatvariables}=1} \, .
    \end{align*}
\end{theorem}
\begin{proof}
    First we realize that the characterization of by $\basemeasure$ representable distributions is a standard simplex extended by trivial coordinates, that is
    \begin{align*}
        \bmrealprobof{\basemeasure}
        = \convhullof{\onehotmapofat{\shortcatindices}{\shortcatvariables}\wcols\basemeasureat{\indexedshortcatvariables}=1 } \, .
    \end{align*}
    This follows from the fact, that the support of any by $\basemeasure$ representable distribution is contained in the support of $\basemeasure$.
    Further, each representable distribution is contained in the convex hull of the one-hot encoded support elements, since any distribution is normalized.

    The polytope of mean parameters is a linear transform of the elements in $\bmrealprobof{\basemeasure}$, since the contraction with $\sencsstat$ is linear.
    It follows that
    \begin{align*}
        \genmeanset
        &= \convhullof{\contractionof{\sencsstatat{\shortcatvariables,\selvariable},\onehotmapofat{\shortcatindices}{\shortcatvariables}}{\selvariable} \wcols\basemeasureat{\indexedshortcatvariables}=1 } \\
        &= \convhullof{\sencsstatat{\indexedshortcatvariables,\selvariable} \wcols\basemeasureat{\indexedshortcatvariables}=1} \, . \qedhere
    \end{align*}
\end{proof}

% Statistic encoding, universal statistic example,
We thus understand the map
\begin{align*}
    \sstatencoding \defcols \facstates\rightarrow\genmeanset \defspace \catindex \rightarrow \sencsstatat{\indexedshortcatvariables,\selvariable}
\end{align*}
as an encoding of states with respect to a statistic $\sstat$.
When the statistic is the universal statistic $\sstat=\universalstat$, this statistic encoding coincides with the one-hot encoding and the mean polytope is the simplex of dimension $\contraction{\basemeasure}-1$ %$(\prod_{\catenumeratorin}\catdimof{\catenumerator})-1$
\begin{align*}
    \meansetof{\universalstat,\basemeasure}
    = \convhullof{\onehotmapofat{\shortcatindices}{\shortcatvariables}\wcols\basemeasureat{\indexedshortcatvariables}=1}\, .
\end{align*}
A generic mean polytope $\genmeanset$ is then the linear transform of the simplex by the contraction with $\sencsstatwith$, where $\selvariable$ is left open.%statistic encoding $\sstatencoding$.


\subsect{Representation as Intersecting Half-Spaces}

For any vector $\normalvec[\selvariable]\in\parspace$ and a scalar $\normalbound\in\rr$, we call the set
\begin{align*}
    \left\{\meanparamwith \wcols\contraction{\meanparamwith,\normalvec[\selvariable]} \leq \normalbound \right\} \subset \parspace
\end{align*}
a half-space of $\rr^\seldim$.
Bounded intersections of finitely many half-spaces are called $\mathcal{H}$-polytopes \cite{ziegler_lectures_2013}.
% Halfspace Representation
We state next, that the polytope $\genmeanset$ of mean parameters is a $\mathcal{H}$-polytopes.


\begin{theorem}
    \label{the:meanPolytopeHalfspaces}
    The set $\genmeanset$ is for any statistic $\sstat$ and base measure $\basemeasure$ a $\mathcal{H}$-polytope, i.e. there exists a finite collection
    \begin{align*}
        \halfspaceparams
    \end{align*}
    where $a_i[\selvariable]$ a vector and $b_i\in\rr$ for all $i\in[n]$ such that
    \begin{align*}
        \genmeanset
        = \left\{\meanparamwith \wcols\forall_{i\in[n]} \, \contraction{\meanparamwith,\normalvecofat{i}{\selvariable}}\leq\normalboundof{i} \right\} \, .
    \end{align*}
\end{theorem}
\begin{proof}
    By \theref{the:meanPolytopeConvHull}, the set $\genmeanset$ is the convex hull of a finite set of vectors and is therefore a $\mathcal{V}$-polytope.
    We therefore apply the main theorem for polytopes \cite{motzkin_beitrage_1936}, which states the equivalence of $\mathcal{V}$-polytopes and $\mathcal{H}$-polytope, for which a proof can be found as Theorem~1.1 in \cite{ziegler_lectures_2013}.
    Therefore, $\genmeanset$ is also a $\mathcal{H}$-polytope and has is thus the intersection of finitely many half-spaces.
\end{proof}

The determination of the half-space parametrizing $\halfspaceparams$ is, however, in general difficult and the main reason for the intractability of probabilistic inference (see e.g. \cite{wainwright_graphical_2008}).


\subsect{Characterization of the Interior}

The interior of the mean polytope consists of the mean parameters to positive distributions as we show next.

\begin{theorem}
    \label{the:meanPolytopeInteriorCharacterization}
    For any minimal statistics $\sstat$ with respect to a boolean base measure $\basemeasure$ (see \defref{def:minimalStatistics}) and a with respect to $\basemeasure$ positive distribution $\probwith$ we have
%    we have for some $\meanparamwith$ that $\meanparamwith\in\sbinteriorof{\genmeanset}$ if and only if there is a positive distribution with respect to $\basemeasure$ such that
    \begin{align*}
        %\meanparamwith
        %=
        \contractionof{\probwith,\sencsstatwith}{\selvariable} \in \sbinteriorof{\genmeanset} \, .
    \end{align*}
\end{theorem}
\begin{proof}
%    \proofrightsymbol: % LATER!
%    By \theref{the:meanPolytopeInterior} we find a canonical parameter $\canparamat{\selvariable}$ such that
%    \begin{align*}
%        \meanparamwith
%        = \contractionof{\expdistat{\shortcatvariables},\sencsstatwith}{\selvariable} \, .
%    \end{align*}
%    We notice that $\expdist$ is positive with respect to $\basemeasure$, as is any member of an exponential family with base measure $\basemeasure$.

    % \proofleftsymbol: % Orient on proof of Theorem~3.3
    Since by assumption the statistics is minimal, the convex set $\genmeanset$ is full dimensional (see e.g. Appendix B in \cite{wainwright_graphical_2008}).
    We thus use a well-known property for full-dimensional convex sets (see \cite{rockafellar_convex_1997,hiriart-urruty_convex_1993}), that $\meanparam\in\interiorof{\genmeanset}$ if for any non-vanishing vector $\vectorat{\selvariable}$ there is a  % citations from Wainwright - Appendix B
    there is a $\tilde{\meanparam}[\selvariable]$ with
    \[ \contraction{\vectorat{\selvariable},\meanparamwith} <  \contraction{\vectorat{\selvariable},\tilde{\meanparam}[\selvariable]} \, . \]
    It thus suffices to show for an arbitrary non-vanishing vector $\vectorat{\selvariable}$ the existence of a distribution $\tilde{\probtensor}$, such that
    \begin{align*}
        \contraction{\vectorat{\selvariable},\meanparamwith} < \contraction{\vectorat{\selvariable},\sencsstatwith,\secprobat{\shortcatvariables}} \, .
    \end{align*}
    We define for $\epsilon\in\rr$
    \begin{align*}
        \probofat{\epsilon}{\shortcatvariables}
        = \normalizationof{\probat{\shortcatvariables},\expof{\epsilon\cdot\contractionof{\sencsstatwith,\vectorat{\selvariable}}{\shortcatvariables}}}{\shortcatvariables}
    \end{align*}
    The derivation of this map at $\epsilon=0$ is
    \begin{align*}
        \difwrt{\epsilon}\probofat{\epsilon}{\shortcatvariables}|_{\epsilon=0}
        = \contractionof{\probwith,\sencsstatwith,\vectorat{\selvariable}}{\shortcatvariables} - \contraction{\probwith,\sencsstatwith,\vectorat{\selvariable}} \cdot \probwith
    \end{align*}
    and thus
    \begin{align*}
        \difwrt{\epsilon} \contraction{\probofat{\epsilon}{\shortcatvariables},\sencsstatwith,\vectorat{\selvariable}}|_{\epsilon=0}
        &= \contractionof{\probwith,(\contractionof{\sencsstatwith,\vectorat{\selvariable}})^2}{\shortcatvariables} \\
        & \quad - \left(\contractionof{\probwith,\sencsstatwith,\vectorat{\selvariable}}{\shortcatvariables}\right)^2 \, .
    \end{align*}
    We can interpret this quantity as the variance of the random variable
    \begin{align*}
        \contractionof{\sencsstatwith,\vectorat{\selvariable}}{\indexedshortcatvariables} \, ,
    \end{align*}
    where $\shortcatindices$ is drawn from $\probwith$.
    The variance is greater than zero, if this random variable is not constant.
    But from the minimality of $\sstat$ with respect to $\basemeasure$ it follows, that this variable is not constant and we therefore have
    \begin{align*}
        0 < \difwrt{\epsilon} \contraction{\probofat{\epsilon}{\shortcatvariables},\sencsstatwith,\vectorat{\selvariable}}|_{\epsilon=0} \, .
    \end{align*}
    Thus, there is a $\epsilon>0$ with
    \begin{align*}
        \contraction{\vectorat{\selvariable},\meanparamwith} < \contraction{\vectorat{\selvariable},\sencsstatwith,\probofat{\epsilon}{\shortcatvariables}} \, . & \qedhere
    \end{align*}
\end{proof}

%
While \theref{the:meanPolytopeInteriorCharacterization} only states that the mean parameter of each positive distribution is in the interior, we will construct for each interior point a positive distribution in \charef{cha:probReasoning} by a member of the corresponding exponential family.

\subsect{Characterization of the Boundary by Faces}

Let us now continue with the investigation of the faces of the mean parameter polytope.%, which we define analogously to Def.~2.1 in \cite{ziegler_lectures_2013}. % ! Defined in Ziegler directly with normals

\begin{definition}
    \label{def:meanPolytopeFaces}
    Given a mean parameter polytope $\genmeanset$ in the half space representation of \theref{the:meanPolytopeHalfspaces}, and any subset $\mathcal{I}\subset[n]$ we say that the set
    \begin{align*}
        \genfacesetof{\facecondset}
        = \left\{\meanparamwith\in\genmeanset \wcols \forall_{i\in\mathcal{I}} \, \contraction{\meanparamwith,\normalvecofat{i}{\selvariable}}=\normalboundof{i} \right\}
    \end{align*}
    is the face to the constraints $\mathcal{I}$.
\end{definition}

While all inequalities in a half-space representation are satisfied for any element of the polytope, we defined faces by the additional sharp satisfaction of a subset of the half-space inequalities.
In this way, the faces build the boundary of $\genmeanset$.
This can be easily verified, since for any vector $\meanparamwith\in\genmeanset$, for which no halfspace inequalities hold sharply, also a neighborhood satisfies the halfspace inequalities.
If any halfspace inequality holds sharply, in the other case, the vector is a member of the corresponding face.

% Trivial face containing the whole polytope in case of non-minimal statistics
If $\sstat$ is not minimal with respect to $\basemeasure$, we find a non-vanishing vector $\vectorat{\selvariable}$ and a scalar $\lambda\in\rr$ such that
\begin{align*}
    \contractionof{\sencsstatat{\shortcatvariables,\selvariable},\vectorat{\selvariable},\basemeasurewith}{\shortcatvariables} = \lambda\cdot\basemeasurewith \, .
\end{align*}
This implies, that any probability distribution $\probwith$ representable with $\basemeasure$ satisfies
\begin{align*}
    \contraction{\probwith,\sencsstatat{\shortcatvariables,\selvariable},\vectorat{\selvariable},\basemeasurewith} = \lambda\cdot \contraction{\probwith,\basemeasurewith} = \lambda \, .
\end{align*}
Any $\meanparamwith\in\genmeanset$ then satisfies
\begin{align*}
    \contraction{\meanparamwith,\vectorat{\selvariable}} = \lambda \, .
\end{align*}
Thus, the polytope $\genmeanset$ is contained in an affine linear subspace and has vanishing interior.
We can further understand this equation as two half-space inequalities
\begin{align*}
    \contraction{\meanparamwith,\vectorat{\selvariable}} \leq \lambda \andspace \contraction{\meanparamwith,\vectorat{\selvariable}} \geq \lambda \, ,
\end{align*}
which can be integrated into any half-space representation.
We conclude, that in the case of non-minimal statistics, the whole polytope $\genmeanset$ is a face itself, since it satisfies these half-space inequalities sharply.


%\subsect{Base measures on faces}

\begin{lemma}
    \label{lem:faceConvHullPreimage}
    For each face $\genfacesetof{\facecondset}$ we have
    \begin{align*}
        \genfacesetof{\facecondset}
        = \convhullof{\sencsstatat{\indexedshortcatvariables,\selvariable}\wcols\shortcatindices\in(\sstatencoding)^{-1}(\genfacesetof{\facecondset})\ncond\basemeasureat{\indexedshortcatvariables}=1} \, .
    \end{align*}
\end{lemma}
\begin{proof}
    This holds, since each face is the convex hull of the contained vertices (see Proposition~2.2 and 2.3 in \cite{ziegler_lectures_2013}).
    Since the vertices are contained in the image of the statistic encoding $\sstatencoding$, the vertices contained in $\genfacesetof{\facecondset}$ are contained in the set
    \begin{align*}
        \sencsstatat{\indexedshortcatvariables,\selvariable}\wcols\shortcatindices\in(\sstatencoding)^{-1}(\genfacesetof{\facecondset}) \, . & \qedhere
    \end{align*}
\end{proof}

\lemref{lem:faceConvHullPreimage} implies in particular, that faces are mean parameter polytopes with respect to refined base measures.
For reference in later chapters, we define these refined base measures next as face measures.

\begin{definition}
    \label{def:faceMeasure}
    The face measure to the face $\genfacesetof{\facecondset}$ of $\genmeanset$ is the boolean tensor $\basemeasureofat{\sstat,\facecondset}{\shortcatvariables}$ with coordinates to $\shortcatindicesin$ by
    \begin{align*}
        \basemeasureofat{\sstat,\facecondset}{\indexedshortcatvariables}
        = \begin{cases}
              1 & \ifspace \sstatencodingof{\shortcatindices}\in\genfacesetof{\facecondset} \\
              0 & \text{else}
        \end{cases} \, .
%        = \indicatorofat{\sstatencodingof{\shortcatindices}\in\genfacesetof{\facecondset}}{\shortcatvariables} \, .
    \end{align*}
\end{definition}

We now specify the mean parameter polytope to any face using the face measure as a refinement of the base measure.

%
\begin{theorem}
    \label{the:faceAsRefinedPolytope}
    For any face $\genfacesetof{\facecondset}$ of $\genmeanset$, we have with the refined base measure
    \begin{align*}
        \secbasemeasureat{\shortcatvariables} = \contractionof{\basemeasureat{\shortcatvariables},\basemeasureofat{\sstat,\facecondset}{\shortcatvariables}}{\shortcatvariables}
    \end{align*}
    that
    \begin{align*}
        \genfacesetof{\facecondset} = \meansetof{\sstat,\secbasemeasure} \, .
    \end{align*}
\end{theorem}
\begin{proof}
    We notice that for any $\shortcatindices\in\facstates$, $\shortcatindices\in(\sstatencoding)^{-1}(\genfacesetof{\facecondset})$ is equal to $\basemeasureofat{\sstat,\facecondset}{\indexedshortcatvariables}=1$ and thus
    \begin{align*}
        \left\{\shortcatindices\wcols\shortcatindices\in(\sstatencoding)^{-1}(\genfacesetof{\facecondset})\ncond\basemeasureat{\indexedshortcatvariables}=1\right\}
        = \left\{\shortcatindices\wcols\secbasemeasureat{\indexedshortcatvariables}=1\right\} \, .
    \end{align*}
    In combination with \lemref{lem:faceConvHullPreimage} we then get
    \begin{align*}
        \genfacesetof{\facecondset}
        &= \convhullof{\sencsstatat{\indexedshortcatvariables,\selvariable}\wcols\shortcatindices\in(\sstatencoding)^{-1}(\genfacesetof{\facecondset})\ncond\basemeasureat{\indexedshortcatvariables}=1} \\
        &= \convhullof{\sencsstatat{\indexedshortcatvariables,\selvariable}\wcols\shortcatindices\wcols\secbasemeasureat{\indexedshortcatvariables}=1} \\
        &= \meansetof{\sstat,\secbasemeasure}
        \, . \qedhere
    \end{align*}
\end{proof}

Positivity of a distribution with respect to face measures is an equivalent condition for the mean parameter of a distribution to be on a face, as we show next.

\begin{theorem}
    \label{the:facePolytopeCharacterization}
    If and only if for a distribution $\probwith$ and a face $\facecondset$ we have
    \begin{align*}
        \contractionof{\probwith,\sencsstatwith}{\selvariable}\in\genfacesetof{\facecondset}\, ,
    \end{align*}
    then $\probwith$ is representable with respect to the base measure
    \begin{align*}
        \contractionof{\basemeasurewith,\genfacemeasurewith}{\shortcatvariables} \, .
    \end{align*}
\end{theorem}
\begin{proof}
    We have
    \begin{align*}
        \meanparamat{\selvariable} = \sum_{\shortcatindices} \probat{\indexedshortcatvariables}\cdot\genstatshortcatencoding \, .
    \end{align*}
    Now, the $\shortcatindices$ with $\genfacemeasureat{\indexedshortcatvariables}=1$ are exactly those, for which the conditions $\facecondset$ hold straight.
    If and only if for a $\shortcatindices$ with $\genfacemeasureat{\indexedshortcatvariables}=0$ we have $\probat{\indexedshortcatvariables}>0$, one of the conditions $\facecondset$ would not hold straight.
    Thus, if and only if $\probwith$ is representable with respect to $\genfacemeasureat{\shortcatvariables}$, we have $\meanparamat{\selvariable}\in\genfacesetof{\facecondset}$.
\end{proof}

% Stronger for exponential distributions: But need to use smallest face for that! -> Maybe later?
For members of exponential families, we can make a stronger statement than \theref{the:facePolytopeCharacterization}.
If for any $\probwith\in\expfamilyof{\sstat,\basemeasure}$ and a face $\facecondset$ we have $\contractionof{\probwith,\sencsstatwith}{\selvariable}\in\sbinteriorof{\genfacesetof{\facecondset}}$ then $\probwith$ is positive with respect to the base measure
\begin{align*}
    \contractionof{\basemeasurewith,\genfacemeasurewith}{\shortcatvariables} \, .
\end{align*}

Let us now investigate tensor network representations of face measures, based on the basis encoding $\bencodingof{\sstat}$ of a statistic.
% Vertices
Vertices of $\genmeanset$ are faces with single elements, that is $\{\meanparamwith\}$.
By \lemref{lem:faceConvHullPreimage} there must be $\meanparam$ must lie in the image of $\sstatencoding$, since otherwise $\genmeanset$ would be empty.
The vertex measure is then
\begin{align*}
    \basemeasureofat{\sstat,\facecondset}{\shortcatvariables}
    = \contractionof{\bencodingofat{\sstat}{\headvariables,\shortcatvariables},\onehotmapofat{\meanparam}{\headvariables}}{\shortcatvariables}
\end{align*}
Here we use that each $\meanparam\in\genfacesetof{\facecondset}\cap\imageof{\sstatencoding}$ has integer-valued coordinates and denote
\begin{align*}
    \onehotmapofat{\meanparam}{\headvariables} = \bigotimes_{\selindexin} \onehotmapofat{\meanparamat{\indexedselvariable}}{\headvariableof{\selindex}} \, .
\end{align*}

\begin{theorem}[Face measure representation]
    \label{the:faceMeasureCharacterization}
    For any face $\genfacesetof{\facecondset}$ of $\meanset$ we have
    \begin{align*}
        \genfacemeasureat{\shortcatvariables}
        =\contractionof{\sstatccwith,\kcoreofat{\facecondset}{\headvariables}}{\shortcatvariables}
    \end{align*}
    where
    \begin{align*}
        \kcoreofat{\facecondset}{\headvariables}
        = \sum_{\meanparam\in\genfacesetof{\facecondset}\cap\imageof{\sstatencoding}} \onehotmapofat{\meanparam}{\headvariables} \, .
    \end{align*}
\end{theorem}
\begin{proof}
    For any $\meanparam\in\genfacesetof{\facecondset}\cap\imageof{\sstatencoding}$ the tensor
    \begin{align*}
        \hypercoreofat{\meanparam}{\shortcatvariables}
        = \contractionof{\sstatccwith,\onehotmapofat{\meanparam}{\headvariables}}{\shortcatvariables}
    \end{align*}
    is the indicator of the preimage of $\meanparam$ under $\sstatencoding$.
    Since preimages the elements in $\genfacesetof{\facecondset}\cap\imageof{\sstatencoding}$ are disjoint, the support of $\hypercoreofat{\meanparam}{\shortcatvariables}$ is disjoint and their sum
    \begin{align*}
        \sum_{\meanparam\in\genfacesetof{\facecondset}\cap\imageof{\sstatencoding}} \hypercoreofat{\meanparam}{\shortcatvariables}
    \end{align*}
    is the indicator of the preimage of $\genfacesetof{\facecondset}$ under $\sstatencoding$, which is the face measure $\basemeasureofat{\sstat,\facecondset}{\shortcatvariables}$.
    Exploiting linearity of contraction we have
    \begin{align*}
        \genfacemeasureat{\shortcatvariables}
        &= \sum_{\meanparam\in\genfacesetof{\facecondset}\cap\imageof{\sstatencoding}} \hypercoreofat{\meanparam}{\shortcatvariables} \\
        &= \contractionof{\bencodingofat{\sstat}{\headvariables,\shortcatvariables},\sum_{\meanparam\in\genfacesetof{\facecondset}\cap\imageof{\sstatencoding}}\onehotmapofat{\meanparam}{\headvariables}}{\shortcatvariables} \\
        &= \contractionof{\bencodingofat{\sstat}{\headvariables,\shortcatvariables},\kcoreofat{\facecondset}{\headvariables}}{\shortcatvariables} \, . \qedhere
    \end{align*}
\end{proof}

%
Let us now investigate, which normalized face measures can be computed using $\sstat$ and a hypergraph $\graph$.

\begin{example}[Vertices]
    \label{exa:vertexMeasures}
    Vertices $\genfaceset$ are proper faces of affine dimension $0$, that is they consist in single vectors.
    Since all vertices are in the image $\sstatencodingat{\stateset}$, there exists an index tuple $\shortcatindices\in\stateset$ such that $\basemeasureat{\indexedshortcatvariables}=1$ and
    \begin{align*}
        \genfaceset = \{\sencsstatat{\indexedshortcatvariables,\selvariable}\} \, .
    \end{align*}
    Then $\kcoreofat{\facecondset}{\headvariables}$ is the one-hot encoding of the by an interpretation map $\indexinterpretation$ assigned index to $\sencsstatat{\indexedshortcatvariables,\selvariable}$, that is
    \begin{align*}
        \kcoreofat{\facecondset}{\headvariables} = \onehotmapofat{\invindexinterpretationat{\sencsstatat{\indexedshortcatvariables,\selvariable}}}{\headvariables} \, .
    \end{align*}
    In particular, the activation core is elementary and the face measure to any vertex is in $\realizabledistsof{\sstat,\elgraph,\basemeasure}$.
\end{example}


% basis CP bound
Extending \exaref{exa:vertexMeasures}, we can provide a coarse estimation of the hypergraph $\graph$ required to decompose $\kcoreof{\facecondset}$ for generic faces $\genfaceset$.
We notice that $\kcoreofat{\facecondset}{\headvariables}$ in \theref{the:faceMeasureCharacterization} is a sparse tensor with basis $\cpformat$ rank $\cardof{\genfacesetof{\facecondset}\cap\imageof{\sstatencoding}}$ (see \charef{cha:sparseRepresentation}).
\begin{align*}
    \sparsityof{\kcoreofat{\facecondset}{\headvariables}}
    = \cardof{
        \sstatencodingat{\genfacemeasure}
    }
\end{align*}
where $\cardof{\sstatencodingat{\genfacemeasure}}$ is the number of different statistic encoding vectors to the support of the face measure $\genfacemeasure$.
Using the formalism of sparse $\cpformat$ decompositions \charef{cha:sparseRepresentation}, this characterize the basis $\cpformat$ rank of $\kcoreofat{\facecondset}{\headvariables}$.
However, the basis $\cpformat$ rank is only an upper bound to generic $\cpformat$ rank, which can be loose.
By \exaref{exa:maximalFaceMeasure} we provide with the maximal face an example, where the basis $\cpformat$ rank is given by the tensor space dimension, whereas the generic $\cpformat$ rank is one and the normalized face measure is thus still in $\realizabledistsof{\sstat,\elgraph}$.

\begin{example}[Maximal face]
    \label{exa:maximalFaceMeasure}
    The maximal face $\genfacesetof{\varnothing}=\genmeanset$ coincides with the mean polytope itself.
    In this case the corresponding activation tensor to the face measure is trivial, that is
    \begin{align*}
        \kcoreofat{\varnothing}{\headvariables} = \onesat{\headvariables} \, .
    \end{align*}
    $\kcoreof{\varnothing}$ is elementary and the normalized face measure $\basemeasureof{\sstat,\varnothing}$ to the maximal face is in $\realizabledistsof{\sstat,\elgraph}$.
\end{example}


\subsect{Partition into Effective Interiors of Faces}

Let us now introduce effective interiors, which enables us to find disjoint partitions of the mean polytope.

\begin{definition}[Effective Interior]
    \label{def:effectiveInterior}
    Let $\arbset\subset\parspace$ be an arbitrary set and $\mathcal{L}$ the minimal affine subspace of $\parspace$ containing $\arbset$.
    Then the effective interior, denoted $\sbinteriorof{\arbset}$ is the the interior of $\arbset$ in the space $\mathcal{L}$.
\end{definition}

\begin{lemma}
    \label{lem:effectiveInteriorPolytopePartition}
    Any polytope is a disjoint union of the effective interiors of its faces, that is
    \begin{align*}
        \genmeanset = \bigcup_{\facecondset\subset[n]}^{\cdot} \sbinteriorof{\genfaceset} \, .
    \end{align*}
\end{lemma}
\begin{proof}
    For any $\meanparam\in\genmeanset$ we find a face such that $\meanparam\in\genfaceset$.
    If $\meanparam\notin\sbinteriorof{\genfaceset}$, then there is a face $\genfacesetof{\tilde{\facecondset}}\subset\genfaceset$ of smaller affine dimension such that $\meanparam\in\genfacesetof{\facecondset}$.
    When continuing this process we reach a face such that $\meanparam\notin\sbinteriorof{\genfaceset}$, since the faces with affine dimension $0$ are vertices and they coincide with their effective interior because they contain a single vector.
\end{proof}

In this way, we find to each $\meanparam\in\genmeanset$ a unique exponential family with statistics $\sstat$ and base measure by a face measure, such that $\meanparam$ is reproduced by an element of that exponential family.
We will show in \charef{cha:probReasoning}, that these reproducing distributions maximize the entropy among any other reproducing distribution.


\subsect{Centers}

We provide further intuition on the position of the mean parameter, when interpreting $\basemeasure$ and $\sstat$ as soft and hard constraint mechanisms on a distribution.
To this end, let $\stateset$ be an arbitrary finite state set, where we so far had the set $\facstates$ for factored representations.

\begin{definition}
    We call the mean parameter of a normalized base measure, that is the vector
    \begin{align*}
        \meanparamofat{\sstat,\zeros,\basemeasure}{\selvariable}
        &= \contractionof{\normalizationof{\basemeasure}{\shortcatvariables},\sencsstatat{\shortcatvariables,\selvariable}}{\selvariable} \, .
    \end{align*}
    the center of $\genmeanset$ (see \figref{fig:hs_mean_sketch}).
\end{definition}

We have
\begin{align*}
    \meanparamofat{\sstat,\zeros,\basemeasure}{\selvariable} = \frac{1}{\contraction{\basemeasure}} \sum_{\shortcatindices\wcols\basemeasureat{\indexedshortcatvariables}=1}  \sencsstatat{\indexedshortcatvariables,\selvariable}
\end{align*}
and therefore understand $\meanparamof{\sstat,\zeros,\basemeasure}$ as the weighted center of $\genmeanset$, where weights are allocated on the image of the statistics encoding by the cardinality of their pre-images.
These concepts are sketched in \figref{fig:hs_mean_sketch} in blue.

%% Mean parameter of base measure
%The mean parameter of a base measure $\basemeasure$, which normalization is understood as a distribution is then
%\begin{align*}
%    \meanparamat{\sstat,\zeros,\basemeasure}
%    &= \frac{1}{\contraction{\basemeasure}} \sum_{\shortcatindices\wcols\basemeasureat{\indexedshortcatvariables}=1}  \sencsstatat{\indexedshortcatvariables,\selvariable} \\
%    &= \contractionof{\normalizationof{\basemeasure}{\shortcatvariables},\sencsstatat{\shortcatvariables,\selvariable}}{\selvariable} \, .
%\end{align*}
%We understand $\meanparamat{\sstat,\zeros,\basemeasure}$ as the weighted center of $\meansetof{\sstat,\basemeasure}$, where weights are allocated on the image of the statistics encoding by the cardinality of their pre-images.
%These concepts are sketched in \figref{fig:hs_mean_sketch} in blue.

The choice of a base measure $\basemeasure$ can be regarded as the hard structure of a distribution.
Given $\basemeasure$, we further implement soft structure in form of distributions in exponential families $\expfamily$.
While the normalized base measure $\basemeasure$ reproduces only the weighted center of $\meansetof{\sstat,\basemeasure}$, canonical parameters can be chosen to alter to any interior point in $\meansetof{\sstat,\basemeasure}$. % ! Here the effective interior, when not minimal !
This alternations are sketched in \figref{fig:hs_mean_sketch} in red.

\begin{figure}
    \begin{center}
        \input{tikz-pics/hs_meanset_sketch}
    \end{center}
    \caption{
        Sketch of hard constraints (blue) and soft constraints (red) determining the position of the mean parameter in $\meanparamof{\sstat,\trivbm}$.
        In the set of states $\stateset$, the support of a base measure $\basemeasure$ is marked blue, and the base measure is the corresponding subset encoding.
        The normalized $\basemeasure$ has the mean parameter $\meanparamof{\sstat,\zeros,\basemeasure}$, and the mean polytope $\meansetof{\sstat,\basemeasure}$ is the convex hull of the image to the statistics encoding $\sstatencoding$ restricted on the support set of $\basemeasure$.
        Given a base measure $\basemeasure$, the corresponding forward and backward maps implement soft contraints depending on canonical parameters $\canparam$.
        The forward mapping maps the zero canonical parameter to $\meanparamof{\sstat,\zeros,\basemeasure}$, correspond with the normalized base measure.
        Generic parameters $\canparam$ (red) are mapped onto the interior $\interiorof{\meansetof{\sstat,\basemeasure}}$.
    }\label{fig:hs_mean_sketch}
\end{figure}